\scanpage{36}
\mspage{32}
% archive sheet number: 35
% source lines 1--3
Il est clair que pour cet $M_c$, les conditions Mag cpt 1 a) b) c) sont satisfaites (pour a)\inserted{ et c)}, on prendra $F=\widetilde X$).

% source lines 4--10
\deleted{\unsure{A}} Si $\unsure{\mathcal F|\mathscr F|\mathfrak F}$ est formé des figures finies, ledit $M_c$ est \unsure{ainsi|alors} la plus grande des compactologies sur $M$. Mais il ne semble pas raisonnable de \uline{définir} en général, $M_c$ en termes du seul \unsure{magasin} $(M,\trianglelefteq,\overset{\circ}{\ll},\mathrel{\mid\circ\mid})$, par les conditions précédentes, comme \deleted{\unsure{les}}
\[
M_c=\{X\in M\mid\forall Y\in M,\ \text{l'ens }\{Y'\in\widetilde Y\mid Y'\mathrel{\overline{\parallel}}X\}\ \text{fini}\},
\]
% source lines 12--21
laquelle donnerait $M_c=M$ quand toutes les multistrates \deleted{de} $\widetilde Y$ de $M$ sont finies. P. ex. dans le cas d'un ensemble ordonné $I$, muni d'une seule extrémité ou \unsure{sans} extrémité (\uline{intervalle ordonné} qui n'est pas un segment), il est naturel de considérer comme compactes les seules strates de la forme $a\ (a\in I)\ I_{a,b}\ (a,b\in I)$.

% source lines 22--26
Il semble bien que ce soit la donnée supplémentaire de la compactologie, qui soit une structure souple plus intéressante que celle d'un atelier \unsure{attenant} particulier.
